\addcontentsline{toc}{chapter}{Summary}
\markboth{Summary}{}
\begin{EnglishSummary}

With the increasing application of artificial intelligence in healthcare, medical tabular data have become an important source for clinical prediction tasks. Such datasets typically contain heterogeneous patient-related variables, including demographic characteristics, laboratory measurements, and physiological indicators. However, learning from medical tabular data remains challenging because of limited sample sizes, heterogeneous feature types, class imbalance, and variation in the predictive relevance of individual variables. Recently, retrieval-augmented tabular learning methods such as TabR have demonstrated promising performance by incorporating information from similar training samples into the prediction process. However, the numerical representation used before retrieval does not explicitly model sample-specific variation in feature emphasis. Therefore, this thesis investigates feature-aware numerical representation for retrieval-augmented medical tabular classification. A feature-aware extension of TabR, termed FiTabR, is proposed by introducing a lightweight attention module before the TabR encoder. The proposed module generates sample-specific attention weights for numerical features and applies residual feature modulation while preserving the original retrieval, contextual aggregation, and prediction mechanisms of TabR. In addition, ablation studies are conducted to examine the contribution of residual preservation, attention scaling, and learned feature weighting to the behavior of the proposed architecture. The learned attention distributions are further analyzed using their mean and standard deviation to characterize feature emphasis and its correspondence with established medical domain knowledge. Experimental results across multiple medical tabular datasets show that FiTabR improves TabR under certain dataset conditions, while the magnitude of improvement varies across datasets. The ablation and attention analyses further indicate that the contribution of the proposed architectural components and the learned feature emphasis is dataset dependent. These findings demonstrate the potential of feature-aware numerical representation for retrieval-augmented tabular learning while also identifying its limitations across different medical data characteristics.


\end{EnglishSummary}
